\section*{Chapitre \ref{ch:g6:producing-our-food} --- Produire notre nourriture}

\begin{solution}{exo:g6:producing-our-food:1}
La \omterm{def:g6:producing-our-food:farming}{culture} gère des \omterm{def:g5:food-webs:roles}{producteurs} --- des \omterm{def:g1:plants-around-us:plant}{plantes} cultivées pour leur
\omterm{def:g6:plants-are-producers:matter}{matière organique}. L'\omterm{def:g6:producing-our-food:farming}{élevage} gère des \omterm{def:g5:food-webs:roles}{consommateurs} --- des \omterm{def:g1:animals-around-us:animal}{animaux}
élevés pour leur production~: \omterm{def:g2:animals-grow-up:born}{lait}, viande, \omterm{def:g2:animals-grow-up:born}{œufs}, laine.
\end{solution}

\begin{solution}{exo:g6:producing-our-food:2}
Grain~: la \omterm{def:g2:from-seed-to-plant:seed}{graine} du blé (sa réserve emballée). Pomme de terre~: le
\omterm{ex:g4:plants-without-seeds:bulbtuber}{tubercule} du pied de pomme de terre. Pomme~: le \omterm{prop:g3:flowers-fruits-seeds:transform}{fruit} du pommier.
\omterm{def:g2:animals-grow-up:born}{Lait}~: la production de soin maternel de la vache. \omterm{def:g2:animals-grow-up:born}{Œuf}~: la nurserie
garnie de la poule.
\end{solution}

\begin{solution}{exo:g6:producing-our-food:3}
La transformation d'une matière première par des \omterm{def:g6:producing-our-food:fermentation}{microbes} qui s'en
nourrissent. Le pain (matière première~: la pâte de farine,
ouvrier~: la \omterm{ex:g6:producing-our-food:bread}{levure}) et le yaourt (matière première~: le \omterm{def:g2:animals-grow-up:born}{lait},
ouvriers~: les \omterm{def:g6:producing-our-food:fermentation}{microbes} acidifiants)~; le fromage et le vinaigre
conviennent aussi.
\end{solution}

\begin{solution}{exo:g6:producing-our-food:4}
La \omterm{ex:g6:producing-our-food:bread}{levure} --- un \omterm{def:g6:decomposers-and-soil:decomposers}{champignon} à une seule \omterm{def:g5:first-look-at-cells:cell}{cellule} --- se nourrit des
sucres de la farine et libère un gaz~; la pâte élastique piège le
gaz en bulles et gonfle. La chaleur du four met fin au travail de la
\omterm{ex:g6:producing-our-food:bread}{levure} et fige la mie levée, trous compris.
\end{solution}

\begin{solution}{exo:g6:producing-our-food:5}
L'acide libéré par les \omterm{def:g6:producing-our-food:fermentation}{microbes} qui se nourrissent du sucre du
\omterm{def:g2:animals-grow-up:born}{lait}~: l'acide épaissit et fait prendre le \omterm{def:g2:animals-grow-up:born}{lait}, et le goût frais et
acidulé est cet acide même.
\end{solution}

\begin{solution}{exo:g6:producing-our-food:6}
Saler, refroidir (le réfrigérateur), sécher --- et aussi cuire~:
chacune refuse aux \omterm{def:g6:producing-our-food:fermentation}{microbes} indésirables une condition de travail
(l'\omterm{def:g6:exploring-our-environment:conditions}{humidité}, la chaleur), exactement comme le froid ralentit un tas
de compost.
\end{solution}

\begin{solution}{exo:g6:producing-our-food:7}
L'eau~: l'agriculteur compte sur la pluie (et parfois l'apporte).
La \omterm{def:g6:exploring-our-environment:conditions}{lumière}~: des semis en plein champ, espacés pour que chaque
\omterm{def:g1:plants-around-us:plant}{plante} ait sa part. La chaleur~: un semis calé sur la saison chaude.
L'air~: un \omterm{def:g6:decomposers-and-soil:soil}{sol} découvert et ameubli. Les \omterm{prop:g4:what-plants-need:needs}{sels minéraux}~: le fumier
et le réapprovisionnement après la récolte.
\end{solution}

\begin{solution}{exo:g6:producing-our-food:8}
(1) Organisme et stade~: le \omterm{def:g2:animals-grow-up:born}{lait} d'une vache --- la production d'une
consommatrice. (2) Place sur le cycle~: production d'un
\omterm{def:g5:food-webs:roles}{consommateur}, achevée par des \omterm{def:g6:producing-our-food:fermentation}{microbes}. (3) Matière première~: le
\omterm{def:g2:animals-grow-up:born}{lait}~; les ouvriers l'acidifient et le font prendre en caillé, puis
affinent le fromage pressé pendant des mois --- trous, croûte et
odeur sont leur journal. (4) Piste~: le \omterm{def:g2:animals-grow-up:born}{lait} vient de la vache, la
vache de l'herbe, l'herbe de la \omterm{def:g6:exploring-our-environment:conditions}{lumière}.
\end{solution}

\begin{solution}{exo:g6:producing-our-food:9}
Les \omterm{def:g6:producing-our-food:fermentation}{microbes} sont des ouvriers vivants avec leurs \omterm{def:g6:exploring-our-environment:conditions}{conditions}~: le
froid arrête presque leur alimentation. Au réfrigérateur, cela
protège l'aliment (les gâcheurs restent inactifs)~; dans le tas
d'hiver, cela retarde le compost (les \omterm{def:g6:decomposers-and-soil:decomposers}{décomposeurs} restent
inactifs). Le principe~: \omterm{prop:g6:decomposers-and-soil:decomposition}{décomposition} et altération avancent au
rythme que les \omterm{def:g6:exploring-our-environment:conditions}{conditions} autorisent.
\end{solution}

\begin{solution}{exo:g6:producing-our-food:10}
«~\omterm{def:g1:healthy-habits:germs}{Microbe}~» désigne une foule immense d'\omterm{def:g6:exploring-our-environment:components}{êtres vivants} minuscules et
très divers. Certains, au mauvais endroit, provoquent des maladies
--- ce sont ceux que vise le lavage des mains~; d'autres, choisis et
gérés dans des \omterm{def:g6:exploring-our-environment:conditions}{conditions} contrôlées, font lever le pain et prendre
le \omterm{def:g2:animals-grow-up:born}{lait}. Ennemi ou employé n'est pas la nature du \omterm{def:g1:healthy-habits:germs}{microbe}, mais
l'accord entre l'\omterm{def:g5:biodiversity:species}{espèce}, le lieu et la gestion.
\end{solution}

\begin{solution}{exo:g6:producing-our-food:11}
Le jambon est passé par un cochon~: son poids a coûté plusieurs fois
son poids en aliments, plus la part brûlée par le cochon, la terre
et l'eau --- le péage de l'étage \omterm{def:g5:food-webs:roles}{consommateur}. Le blé du pain
n'était qu'à un étage de la \omterm{def:g6:exploring-our-environment:conditions}{lumière}. Gaspiller le sandwich gaspille
les deux registres~; celui du jambon est le plus long.
\end{solution}

\begin{solution}{exo:g6:producing-our-food:12}
Ils ont géré les \emph{\omterm{def:g6:exploring-our-environment:conditions}{conditions}} des ouvriers~: la chaleur du coin
où repose la pâte, le levain gardé et nourri, le sel et les temps
réglés par tâtonnement au fil des générations. Les recettes
marchaient parce que maintenir de bonnes \omterm{def:g6:exploring-our-environment:conditions}{conditions} fait prospérer
les bons \omterm{def:g6:exploring-our-environment:components}{êtres vivants} --- que l'on sache ou non que ces \omterm{def:g6:exploring-our-environment:components}{êtres
vivants} sont là. Ils ont élevé des \omterm{def:g6:producing-our-food:fermentation}{microbes} à l'aveugle, et bien.
\end{solution}

\begin{solution}{pb:g6:producing-our-food:1}
\textbf{1.} Chaque grain est une \omterm{def:g2:from-seed-to-plant:seed}{graine}~: une minuscule \omterm{def:g1:plants-around-us:plant}{plante} de
blé garnie d'une \omterm{def:g3:animals-through-seasons:active}{réserve de nourriture}, faite par la mère pour les
premiers jours de cette plantule --- pas pour nous, même si nous la
détournons.

\textbf{2.} Dans les \omterm{def:g1:plants-around-us:parts}{feuilles} vertes des pieds de blé eux-mêmes, à
partir de l'eau et des \omterm{prop:g4:what-plants-need:needs}{sels minéraux} du \omterm{def:g6:decomposers-and-soil:soil}{sol} et du dioxyde de carbone
de l'air, avec la \omterm{def:g6:exploring-our-environment:conditions}{lumière} pour \omterm{prop:g3:balanced-diet:jobs}{énergie}.

\textbf{3.} Semé, un grain parcourt son \omterm{def:g2:born-grow-die:cycle}{cycle de vie}~: \omterm{def:g2:from-seed-to-plant:germination}{germination},
puis \omterm{def:g1:plants-around-us:plant}{plante}, puis nouveaux grains. Vendu et moulu, son \omterm{def:g2:born-grow-die:cycle}{cycle de vie}
s'arrête et sa réserve devient la nôtre --- un casse-croûte emballé
pour une plantule, mangé par une personne.

\textbf{4.} Rendre ce qu'on emprunte~: en se décomposant, le fumier
réapprovisionne les \omterm{prop:g4:what-plants-need:needs}{sels minéraux} que la récolte a emportés --- il
sert le besoin en \omterm{prop:g4:what-plants-need:needs}{sels minéraux}.

\textbf{5.} Parce que la \omterm{ex:g6:producing-our-food:bread}{levure} est un ouvrier vivant~: au chaud,
elle se nourrit et travaille vite~; au froid, elle reste inactive
--- la règle de l'équipe du compost, à l'échelle d'une boulangerie.

\textbf{6.} Elle s'est nourrie des sucres de la farine et a libéré
du gaz~; la pâte élastique a piégé le gaz en bulles, et les bulles
ont doublé la pâte.

\textbf{7.} La fenêtre froide a mis la \omterm{ex:g6:producing-our-food:bread}{levure} au repos --- à peine
nourrie, à peine de gaz. Les deux équipes ont mené sans le vouloir
une expérience honnête sur la chaleur~: même pâte, même jour, une
seule condition différente, deux pousses à comparer.

\textbf{8.} La chaleur a tué la \omterm{ex:g6:producing-our-food:bread}{levure} et mis fin au service~; la
mie figée garde les trous des bulles --- la signature de la
main-d'œuvre, cuite sur place.

\textbf{9.} La tranche vient de la miche, la miche de la pâte, la
pâte de la farine, la farine du grain, le grain des \omterm{def:g1:plants-around-us:parts}{feuilles} du blé
--- fabriqué là à partir d'entrées minérales avec la \omterm{def:g6:exploring-our-environment:conditions}{lumière}~: la
piste se termine au soleil du champ de juillet.

\textbf{10.} Tranche~: un étage (du \omterm{def:g5:food-webs:roles}{producteur} à la table, avec une
finition par des \omterm{def:g6:producing-our-food:fermentation}{microbes}). Beurre~: deux étages (herbe, vache,
crème). Jambon~: deux étages avec le plus gros péage --- le cochon a
brûlé et rejeté la plus grande part de la matière de sa ration avant
qu'un peu ne devienne du jambon. La tranche est la moins chère en
matière~; le jambon la plus chère.

\textbf{11.} Les \omterm{def:g6:decomposers-and-soil:decomposers}{décomposeurs} du compost --- vers, cloportes,
\omterm{def:g6:decomposers-and-soil:decomposers}{champignons}, équipes invisibles --- démontent la \omterm{def:g6:plants-are-producers:matter}{matière organique}
de la pâte en \omterm{def:g6:decomposers-and-soil:soil}{humus} et en \omterm{prop:g4:what-plants-need:needs}{sels minéraux}~: un tour de plus, de
l'aliment à la fertilité du \omterm{def:g6:decomposers-and-soil:soil}{sol} de l'année suivante.

\textbf{12.} Par exemple~: «~Notre pain a été bâti par un \omterm{def:g5:food-webs:roles}{producteur}
à partir de \omterm{def:g6:plants-are-producers:matter}{matière minérale} et de \omterm{def:g6:exploring-our-environment:conditions}{lumière}, levé par une
main-d'œuvre de \omterm{def:g6:producing-our-food:fermentation}{microbes}, et n'a eu besoin d'aucun étage de
\omterm{def:g5:food-webs:roles}{consommateur} --- alors que les sandwichs au jambon d'à côté en ont
eu besoin d'un.~»
\end{solution}
